Recent open-source finite-element and PDE-optimization frameworks cover
different parts of the workflow needed for nonlinear solves and design. High-level
finite-element systems such as \fenics{} and \code{DOLFINx} emphasize expressive
variational modeling together with compiled kernels and parallel assembly
\citep{logg2012fenicsbook,baratta2025dolfinx}. PDE-level adjoint frameworks
such as \code{dolfin-adjoint}, \code{pyadjoint}, and \code{cashocs} automate
adjoint-centered optimization workflows and PDE-level sensitivity automation
\citep{farrell2013dolfinadjoint,mitusch2019pyadjoint,blauth2023cashocsv2}.
JAX-native or JAX-connected workflows, including JAX-FEM, AutoPDEx,
JAX-CPFEM, the Firedrake--JAX bridge, and recent \code{FEniCSx} external-operator
work, push automatic differentiation deeper into constitutive modeling and
differentiable simulation
\citep{xue2023jaxfem,bode2025autopdex,hu2025jaxcpfem,yashchuk2023bringing,latyshev2025externaloperators}.
Recent topology-optimization implementations such as \code{FEniTop} show that
compact parallel design-mechanics workflows can also be built on modern
\code{FEniCSx} infrastructure \citep{jia2024fenitop}.

This paper targets a narrower design point than the broad differentiable-PDE
landscape above. The maintained production path in \repo{} is a solver-centric
\jaxpetsc{} stack for distributed nonlinear finite-element solves and optimization.
Local energies and constitutive laws are written in JAX, while sparse assembly,
globalization, Krylov solves, and preconditioning are handled by PETSc. The goal
is not to claim one universally preferred framework, but to document one specific
combination that is practically useful for our benchmark families: high-level
local differentiation coupled to distributed sparse Newton--Krylov execution on
hard nonlinear mechanics and design problems.

The central claim is therefore deliberately limited. We do not claim that every
benchmark requires exact Hessians, nor that mainstream frameworks cannot provide
them. We instead show that several derivative routes can be made operational
inside one maintained distributed workflow:
\begin{enumerate}
  \item exact element-level derivatives when the local state dimension is still
        affordable;
  \item constitutive-level automatic differentiation for high-order plasticity,
        where local tangents are much cheaper than full element Hessians;
  \item colored sparse finite-difference (SFD) recovery when exact second-order
        information is not the right cost-quality tradeoff.
\end{enumerate}

These routes are evaluated on six maintained benchmark families:
$p$-Laplace, scalar Ginzburg--Landau, hyperelasticity, Plasticity2D,
\plasticitythree{}, and topology optimization. The paper gives all six families
top-level space because the main contribution is methodological rather than a
single benchmark-specific solver. At the same time, \plasticitythree{} and
topology optimization receive the most detailed comparison because they are the
closest stress tests for the maintained \jaxpetsc{} path; \plasticitythree{} is
also the family in which source-faithfulness, surrogate interpretation, and
historical scaling have to be separated most carefully. The geomechanics context
for that family follows the finite-element strength-reduction and 3D
slope-stability literature \citep{tschuchnigg2015nonassociated,sysala2021optimization,sysala2025advancedcontinuation},
but the repository object studied here is explicitly the implemented endpoint
surrogate, not a claim of path-consistent incremental-history equivalence.

\subsection{State of the Art and Positioning}
\label{sec:intro-sota}

For the purposes of this paper, the most relevant comparison axes are the
modeling layer, the scope of differentiation, the form in which second-order
information enters the workflow, the parallel execution model, and the benchmark
families closest to our maintained suite. \Cref{tab:sota-comparison} summarizes
that comparison using conservative, source-backed descriptors.

\begin{table}[tbp]
  \centering
  \scriptsize
  \setlength{\tabcolsep}{3pt}
  \begin{tabular*}{\textwidth}{@{}>{\RaggedRight\arraybackslash}p{0.16\textwidth}@{\extracolsep{\fill}} >{\RaggedRight\arraybackslash}p{0.13\textwidth} >{\RaggedRight\arraybackslash}p{0.15\textwidth} >{\RaggedRight\arraybackslash}p{0.15\textwidth} >{\RaggedRight\arraybackslash}p{0.14\textwidth} >{\RaggedRight\arraybackslash}p{0.18\textwidth}@{}}
\toprule
Family & Modeling & Differentiation & Second-order route & Parallel & Closest overlap \\
\midrule
\shortstack[l]{FEniCS\\DOLFINx\\\citep{logg2012fenicsbook,baratta2025dolfinx}} & High-level variational forms & Problem-specific; AD is not the central claim & Manual or application-specific & Distributed FEM assembly and solve & Elliptic and finite-strain mechanics \\
\shortstack[l]{dolfin-adjoint\\pyadjoint\\cashocs\\\citep{farrell2013dolfinadjoint,mitusch2019pyadjoint,blauth2023cashocsv2}} & High-level PDE plus optimization loop & Adjoint-based first-order sensitivities & Not the advertised main path & MPI via the host FEM stack & PDE control, shape, and topology \\
\shortstack[l]{JAX-FEM\\AutoPDEx\\\citep{xue2023jaxfem,bode2025autopdex}} & JAX-native PDE / FE programs & Program-level forward and reverse AD & Higher-order JAX derivatives & JAX execution; PETSc optional in AutoPDEx & Nonlinear mechanics and inverse design \\
\shortstack[l]{Firedrake--JAX\\FEniCSx ext. ops\\\citep{yashchuk2023bringing,latyshev2025externaloperators}} & Host FEM stack plus AD bridge & Tangent, adjoint, or local constitutive AD & Local-operator derivatives, not full sparse Hessians & Host framework parallel back end & Parameterized PDEs and constitutive models \\
\shortstack[l]{JAX-CPFEM\\\citep{hu2025jaxcpfem}} & JAX-native crystal-plasticity FEM & Differentiable constitutive simulator & AD constitutive derivatives & GPU-oriented execution & Crystal plasticity \\
\shortstack[l]{FEniTop\\\citep{jia2024fenitop}} & FEniCSx topology code & Sensitivity-based design updates & Not a second-order solver paper & Parallel FEniCSx workflow & 2D and 3D topology optimization \\
This work & JAX local energies plus PETSc sparse solvers & Element AD, constitutive AD, and colored SFD & Local Hessians/tangents or sparse recovery & PETSc MPI vectors, matrices, SNES, and multigrid & $p$-Laplace, GL, hyperelasticity, plasticity, topology \\
\bottomrule
\end{tabular*}

  \caption{State-of-the-art positioning of the framework families most relevant
  to this paper. The columns record source-supported workflow properties rather
  than subjective scores.}
  \label{tab:sota-comparison}
\end{table}

Three comparisons matter most for the narrative that follows. First, relative to
high-level FEM and adjoint ecosystems
\citep{logg2012fenicsbook,baratta2025dolfinx,farrell2013dolfinadjoint,mitusch2019pyadjoint,blauth2023cashocsv2},
\repo{} moves the center of gravity from PDE-level adjoint-centered automation toward
distributed nonlinear solve policy and multiple second-order routes inside the
same sparse solver stack. Second, relative to JAX-native mechanics and
differentiable-simulation frameworks
\citep{xue2023jaxfem,bode2025autopdex,hu2025jaxcpfem}, the main distinction is
that the maintained implementation here centers PETSc-owned sparse MPI data
structures and preconditioned Newton solves, whereas the cited JAX-native
frameworks emphasize JAX-level modeling, local differentiation, GPU execution,
or implicit-differentiation workflows. Third, relative to recent bridge papers
that combine JAX-style differentiation with established FEM environments
\citep{yashchuk2023bringing,latyshev2025externaloperators}, the paper's emphasis
is on extending that bridge-style differentiation story into benchmarked PETSc
sparse assembly, Krylov solves, and nonlinear globalization for benchmark
families that include topology optimization and 3D Mohr--Coulomb-type
mechanics.

The resulting positioning statement is intentionally modest. \repo{} is not
presented as a general differentiable-PDE platform, nor as a replacement for the
high-level FEM and adjoint ecosystems cited above. It is presented as a
maintained JAX+PETSc workflow for distributed nonlinear energy minimization and
equilibrium solves in which derivative fidelity, sparse linear algebra, and
preconditioner policy are all part of the experimental design and comparison
criteria rather than merely implementation detail.

The paper is organized as follows. \Cref{sec:related} deepens the literature
context behind the summary above. \Cref{sec:methodology} describes the common
optimization and derivative machinery, including line search, trust regions, and
colored derivative recovery. \Cref{sec:implementation} explains how these ideas
are mapped onto \purejax{}, \fenics{}, and \jaxpetsc{}. \Cref{sec:benchmarks}
defines the benchmark suite, with each benchmark separated into literature model
and repository specialization. \Cref{sec:numerical-examples} then reports the
timing, scaling, and discretization-comparison evidence, followed by supporting
solver evidence and a discussion of limitations.
